\section{Разложение двенадцатилучевой звезды кубооктаэдра по плоскому срезу $\{111\}$}

Для любого выбранного плоского среза $\{111\}$ двенадцать ближайших соседей звезды~$N_{12}(x)$ гранецентрированной кубической решётки распадаются ровно на три параллельных слоя $(3 + 6 + 3 = 12)$:
\begin{enumerate}[noitemsep]
\item \textbf{Экваториальный слой в плоскости среза ($|N_6(x)| = 6$ узлов):}
  шесть соседей, лежащих в той же плоскости $\{111\}$, что и центральный узел~$x$, образуют правильный плоский шестиугольник с угловым шагом $\pi/3 = 60^\circ$.
\item \textbf{Верхний и нижний треугольные слои ($3 + 3 = 6$ узлов):}
  три соседа в примыкающем слое~$B$ над плоскостью среза образуют правильный треугольник~$K_3$, а три соседа в примыкающем слое~$C$ под плоскостью среза образуют повёрнутый на~$\pi/3$ треугольник~$K_3$.
\end{enumerate}

Такое разложение звезды $\geomNeighborTwelve = 6_{\mathrm{in\text{-}plane}} + 3_{\mathrm{up}} + 3_{\mathrm{down}}$ обеспечивает строгую согласованность между двумерной динамикой в срезе $\{111\}$ (где циркуляция фазы по шестиугольнику~$N_6$ задаёт элементарный топологический контур вихря) и полным трёхмерным переносом по двенадцати связям кубооктаэдра~\latGCKshort.
